\newpage
\thispagestyle{empty}
%
\begin{spacing}{1.3}
%
\chapter{Use Case Analysis}

This chapter describes the primary use cases of the GramaGIS system. Each use case represents a distinct interaction between the actors and the platform, outlining the functional behavior, the technical workflow involving the frontend, Express.js backend, GeoServer services, and the expected outcomes for village administration \cite{fu2020webgis}.

\section{View Interactive Village Map}

\textbf{Actors:} Public Viewers, Administrative Users

\textbf{Description:}
This use case allows any user to access the GramaGIS web portal to explore the spatial distribution of village assets and ward boundaries.

\textbf{Workflow:}
\begin{itemize}
    \item User navigates to the GramaGIS URL in a web browser.
    \item The frontend loads the map interface and requests the base map and published GIS layers.
    \item GeoServer provides the required WMS layers to the Leaflet frontend \cite{fu2020webgis}.
    \item The interactive map is rendered, allowing the user to pan and zoom across the Panchayat.
\end{itemize}

\textbf{Outcome:}
Users gain immediate spatial awareness of the village layout and can identify the locations of public utilities and administrative boundaries.

\section{Natural Language Infrastructure Query}

\textbf{Actors:} Public Viewers, Administrative Users

\textbf{Description:}
Users interact with the geospatial database using plain English queries to find specific infrastructure instead of using manual filters \cite{liu2025nalspatial}.

\textbf{Workflow:}
\begin{itemize}
    \item User enters a query (e.g., ``Where are the clinics in Ward 3?'') into the search interface.
    \item The Express.js backend forwards the text along with the schema hint to the Gemini API.
    \item Gemini returns a JSON object containing the target layer and a valid CQL filter \cite{liu2025nalspatial}.
    \item The frontend applies the CQL filter to the relevant map layer and requests filtered features through the WFS proxy.
    \item The map and result panel update to highlight only the relevant features.
\end{itemize}

\textbf{Outcome:}
Non-technical users can perform spatial searches without knowledge of SQL or GIS query syntax.

\section{Metadata Retrieval}

\textbf{Actors:} Public Viewers, Administrative Users

\textbf{Description:}
Retrieval of detailed non-spatial information associated with a specific map feature, such as a school or hospital.

\textbf{Workflow:}
\begin{itemize}
    \item User clicks on a specific feature on the interactive map.
    \item Leaflet issues a \texttt{GetFeatureInfo} request to the active GeoServer layer.
    \item GeoServer returns the matching feature attributes from the published spatial data \cite{postgis_in_action}.
    \item A formatted details panel is displayed on the frontend containing the feature metadata.
\end{itemize}

\textbf{Outcome:}
Users access descriptive data linked directly to the geographic location of the asset.

\section{Spatial Data Export}

\textbf{Actors:} Administrative Users

\textbf{Description:}
Exporting filtered layer data for offline use, official reports, or secondary analysis \cite{fu2020webgis}.

\textbf{Workflow:}
\begin{itemize}
    \item User applies filters to the map or selects a target layer in the admin module.
    \item User selects the ``Export'' function and chooses a format exposed by the interface, such as GeoJSON or CSV.
    \item The Express.js proxy forwards the export request to GeoServer.
    \item GeoServer generates the filtered export from the underlying spatial data and returns it through the backend.
\end{itemize}

\textbf{Outcome:}
Administrators can produce standardized data files for documentation and sharing with other departments.

\section{Administrative Data Update (T-GIS)}

\textbf{Actors:} Administrative Users

\textbf{Description:}
Authorized Panchayat officials update or add infrastructure data to the living database to reflect real-world changes \cite{smart_panchayat2023}.

\textbf{Workflow:}
\begin{itemize}
    \item Administrator logs in and receives a JSON Web Token (JWT).
    \item The user selects an existing feature to edit or uses the custom geometry editor to prepare an updated geometry.
    \item The frontend sends the updated attributes, geometry in WKT form, and the JWT to the Express.js secure proxy route.
    \item Express.js validates the token and translates the payload into a WFS-T (Web Feature Service-Transactional) request \cite{fu2020webgis}.
    \item GeoServer commits the change to the PostGIS database \cite{postgis_in_action}.
\end{itemize}

\textbf{Outcome:}
The village database remains accurate and up to date, providing a reliable source of truth for administrative planning.

\section{Infrastructure Resource Review}

\textbf{Actors:} Administrative Users

\textbf{Description:}
Administrators review filtered records and live layer data to inspect the distribution of mapped infrastructure and verify whether datasets are up to date.

\textbf{Workflow:}
\begin{itemize}
    \item Administrator opens the management interface and selects a layer such as schools, hospitals, roads, or government offices.
    \item The interface loads the schema and current feature set through the backend and GeoServer services.
    \item The administrator filters or searches records to inspect coverage, attributes, and current layer contents.
    \item The system presents the matching records in the table and the associated features on the map.
\end{itemize}

\textbf{Outcome:}
Panchayat officials gain a practical operational view of currently published infrastructure records and can identify missing or outdated data for correction.

\section{Feedback and Infrastructure Issue Reporting}

\textbf{Actors:} Public Viewers, Administrative Users

\textbf{Description:}
This use case allows users to report inaccuracies in map data, missing infrastructure elements, or service-related issues directly to the Panchayat office.

\textbf{Workflow:}
\begin{itemize}
    \item User opens the general ``Feedback'' interface.
    \item The user enters the issue details, category, ward, and optional coordinates or landmark information.
    \item The frontend sends a POST request with the feedback payload to the Express.js API.
    \item The backend validates the input and stores the report in the \texttt{feedback} table through direct SQL operations.
    \item An administrative user reviews the report, updates its workflow status, and uses the T-GIS module if a spatial data correction is required \cite{smart_panchayat2023}.
\end{itemize}

\textbf{Outcome:}
GramaGIS maintains a higher level of data accuracy and operational transparency through citizen reporting and continuous administrative oversight.

\end{spacing}
